%!TEX root = ../../main.tex
\section{평가 규칙}
\label{sec:pred-evaluation}

\begin{itemize}[leftmargin=1.5em]
  \item \textbf{점수.} 적정 점수인 Brier 점수와 로그 손실(낮을수록 좋음)을 주로, AUC를 보조로 보고한다.
  \item \textbf{셀.} 모든 평가 셀은 코호트별(T 또는 S)이다. 코호트 사이의 진술은 동일 행에서 예비 선언된 전이·합동 대비(제 \ref{sec:pred-transfer} 절)와 공통 경기 대비(제 \ref{sec:val-external} 절)로 한정한다.
  \item \textbf{가중.} 모든 평가 셀에서 각 교전(또는 시간상태 질의)은 가중치 $w_i = 1/n_m$을 갖는다. $n_m$은 그 경기가 셀 안에 갖는 행 수이며 셀마다 다시 계산한다. Brier, 로그 손실, AUC, 양성 비율, CORP 성분 모두 이 가중 평균이다. 예외는 발생하는 곳에 표기한다(비교전 대조의 중앙값·분위수는 행 단위, 킬 축 불일치 비율과 다음 오브젝트 건수는 비가중).
  \item \textbf{짝 대비와 부트스트랩.} $\dBrier(q - \text{기준선})$은 경기를 복원 추출하고 경기별 가중 제곱오차 합을 다시 모아 \numBootDraws{}회(시드 \bootSeed) 추출의 2.5·97.5 백분위를 95\% 구간으로 삼는 짝 경기 군집 백분위 부트스트랩으로 평가한다. 양수 추출 비율은 $\dBrier > 0$인 추출의 비율이며 $p$값과 다른 양이다. 구간은 평가기·라벨·선정이 고정된 아래의 경기 표본 변동을 기술하며, 모델 적합이나 라벨 모형의 불확실성은 포함하지 않는다. 주 결과의 열네 개 안팎의 구간은 보정 없이 보고하고, 외부 표본의 보완 부트스트랩(제 \ref{sec:val-external} 절)은 네 대비 가족의 Bonferroni 구간을 함께 보고한다.
  \item \textbf{구간이 있는 분석과 없는 분석.} 구간은 짝 $\dBrier$ 대비, 짝지은 비교전 대조, 정보군 대비, 외부 보완 부트스트랩에 있다. $V \to W$ 표, 대응 표, 지평 표, CORP 성분, 층화 표, 점수 전용 외부 표의 값은 점값이며 각 표의 표주에 그렇게 적는다.
  \item \textbf{해석 문턱.} $\dBrier$의 ``의미 있는'' 크기로 $\tau = \tauMeaningful$을 해석용으로 사전 선언하였다. 이 값은 검정 기준과 무관한 해석의 눈금이다.
  \item \textbf{증거의 지위.} 본 논문의 모든 수치는 동결된 실험 원장에서 옮긴 값이며, 두 코호트의 주 대비는 동결 예측표에서 다시 집계되어 문서값과 일치한다(부록 \ref{app:repro}). 승률·시간 기준선을 유연하게 만든 주 대비, 점수 분해, 외부 평가는 초기 분석에서 시험 패치 15.16이 노출된 뒤 수행되었으므로 사전 등록된 확증 검정과 구분하여 \emph{탐색적} 결과로 표기한다. 소규모 교전 코호트 S는 사전 선언된 계약 아래 추가되었고 같은 지위를 갖는다. 보완 실험(제 \ref{sec:eng-endpoint}, \ref{sec:eng-sensitivity}, \ref{sec:value-correspondence}--\ref{sec:value-shold}, \ref{sec:pred-infogroups}, \ref{sec:val-frame}, \ref{sec:val-external} 절)은 15.16 노출 이후의 진단이며 헤드라인을 대체하지 않는다. 실행 상태는 부록 \ref{app:deferred}의 표 \ref{tab:status}에 하위 분석 단위로 적는다.
\end{itemize}
